\documentclass[10pt,a4paper]{amsart}
\usepackage[margin=19mm]{geometry}
\usepackage{amsmath,amssymb,mathtools}
\usepackage[scr=rsfs,cal=euler]{mathalfa}
\usepackage{xfrac}
\usepackage[unicode,hidelinks,bookmarks=false]{hyperref}
\newcommand{\ie}{i.e.\ }
\newcommand{\cf}{cf.\ }
\newcommand{\resp}{respectively\ }
\newcommand{\Subsection}{\subsection}
\providecommand{\nobreakdash}{\nobreak\hbox{-}}
\setlength{\emergencystretch}{2em}
\multlinegap0pt
\usepackage{amssymb}
\usepackage[shortlabels]{enumitem}
\usepackage{tikz}
\usepackage{float}
\usepackage{mathtools}
\newtheorem{proposition}{Proposition}
\newtheorem{theorem}{Theorem}
\newcommand{\dens}{\mathrm{dens}}
\title{Analysis of the Density of Words under Morphism $\{a,b\}$}
\author{Jasem Hamoud, Duaa Abdullah}
\date{Source: arXiv:2601.06150v2 (CC BY 4.0)}
\begin{document}
\maketitle

\noindent\emph{Source and licence.} Analysis of the Density of Words under Morphism $\{a,b\}$. Version arXiv:2601.06150v2 (CC BY 4.0), distributed by arXiv under the Creative Commons Attribution 4.0 International licence, http://creativecommons.org/licenses/by/4.0/, as recorded in the arXiv licence metadata for this exact version and shown on the abstract page https://arxiv.org/abs/2601.06150v2. Full licence notice: \texttt{licenses/arxiv-2601.06150-CC-BY-4.0.txt}.\par\smallskip

\noindent\textit{Editorial scope.} Counted unit: the article's theorem thm1Asymptotic (source0.tex lines 620--634). The article's complete original proof of it is delivered with it (source0.tex lines 636--674, one \texttt{proof} environment of 39 lines). No mathematics is written, repaired or re-derived: the statement and the proof are the article's own lines, in the article's own notation, and no retained line is substituted or re-flowed.  Retained uncounted as the source-local context the proof needs: lines 607--612.  Nothing was omitted from any retained span, so the package carries no omission record.  No retained line contains a citation, so the wrapper carries no bibliography and no bibliography entry.  No retained line contains a figure environment, an image-inclusion command or drawing code, so no image is delivered and the package declares no asset dependency.  Editorial formatting only, in addition: an AMS \texttt{amsart} preamble that restates the article's own declarations and macros used by the retained lines, namely the environments \texttt{proposition}, \texttt{theorem} on separate counters, together with the standard packages those lines need.  The retained environments print in this document as Proposition 1, Theorem 1 in order of appearance, whereas the article prints the same items with the numbering of its own full document.  The complete native source of the article as distributed in the arXiv e-print for this exact version is bundled unchanged as \texttt{sources/192391/source0.tex}.
\begin{proposition}\label{pro1defnewwordn2}
For every $m \ge 1$, the length of $q_m$ is
\[
|q_m| = |y_m| + 2 = F_{m+2} + 2.
\]
\end{proposition}

\begin{theorem}[Asymptotic densities in $\mathbb{Q}$]\label{thm1Asymptotic}
Let $\mathbb{Y} = (y_n)_{n \ge 0}$ and $\mathbb{Q} = (q_m)_{m \ge 1}$ be as above. Define the letter densities by
\[
\dens_a(\mathbb{Q}) := \lim_{m \to \infty} \frac{\#\{\text{$a$'s in } q_m\}}{|q_m|}, 
\qquad
\dens_b(\mathbb{Q}) := \lim_{m \to \infty} \frac{\#\{\text{$b$'s in } q_m\}}{|q_m|},
\]
provided the limits exist. Then
\[
\dens_a(\mathbb{Q}) = \frac{1}{\varphi},
\qquad
\dens_b(\mathbb{Q}) = \frac{1}{\varphi^2},
\]
where $\varphi = \dfrac{1+\sqrt{5}}{2}$ is the golden ratio.
\end{theorem}

\begin{proof}
Let $A_n$ be the number of occurrences of $a$ in $y_n$ and $B_n$ the number of occurrences of $b$ in $y_n$. From the recursion $y_n = y_{n-1}y_{n-2}$ we obtain
\[
A_n = A_{n-1} + A_{n-2},\qquad B_n = B_{n-1} + B_{n-2} \qquad (n \ge 2),
\]
with initial conditions $(A_0,B_0) = (1,0)$ and $(A_1,B_1) = (1,1)$. Solving these recurrences yields
\[
A_n = F_{n+1},\qquad B_n = F_n \qquad (n \ge 0),
\]
where $(F_n)_{n\ge 0}$ is the Fibonacci sequence with $F_0=0$, $F_1=1$.

For the framed word $q_m = a y_m b$, we therefore have
\[
\#\{\text{$a$'s in } q_m\} = A_m + 1 = F_{m+1} + 1,\qquad
\#\{\text{$b$'s in } q_m\} = B_m + 1 = F_m + 1,
\]
and by Proposition~\ref{pro1defnewwordn2},
\[
|q_m| = |y_m| + 2 = F_{m+2} + 2.
\]
Hence
\begin{equation}\label{eq2thm1Asymptotic}
\dens_a(\mathbb{Q}) = \lim_{m \to \infty} \frac{F_{m+1} + 1}{F_{m+2} + 2},
\qquad
\dens_b(\mathbb{Q}) = \lim_{m \to \infty} \frac{F_m + 1}{F_{m+2} + 2}.
\end{equation}
It is well known that
\[
\lim_{k \to \infty} \frac{F_k}{F_{k+1}} = \frac{1}{\varphi},
\qquad
\lim_{k \to \infty} \frac{F_k}{F_{k+2}} = \frac{1}{\varphi^2},
\]
and adding fixed constants to numerator and denominator does not change these limits. Applying these facts to~\eqref{eq2thm1Asymptotic} yields
\[
\dens_a(\mathbb{Q}) = \frac{1}{\varphi},\qquad
\dens_b(\mathbb{Q}) = \frac{1}{\varphi^2},
\]
as claimed.
\end{proof}

\end{document}
